\documentclass[11pt,a4paper]{article}
%-------------------------
% Basic packages
%-------------------------
\usepackage[utf8]{inputenc}
\usepackage[T1]{fontenc}
\usepackage{lmodern}
\usepackage{microtype}
\usepackage{geometry}
\geometry{margin=1in}

% Mathematics
\usepackage{amsmath,amssymb,amsthm,mathtools,bm}
\usepackage{physics}    % convenient bra-ket, derivatives, etc.
\usepackage{siunitx}    % units
\usepackage{braket}     % alternate bra/ket macros
\usepackage{mhchem}     % chemistry notation if needed

% Figures & graphics
\usepackage{graphicx}
\usepackage{tikz}
\usepackage{pgfplots}
\pgfplotsset{compat=1.18}

% References, hyperlinks, and clever names
\usepackage[colorlinks=true,linkcolor=blue,citecolor=blue,urlcolor=blue]{hyperref}
\usepackage[nameinlink]{cleveref}

% Lists and layout
\usepackage{enumitem}
\setlist{itemsep=2pt,parsep=2pt}

% Theorem environments
\theoremstyle{plain}
\newtheorem{theorem}{Theorem}[section]
\newtheorem{lemma}[theorem]{Lemma}
\newtheorem{proposition}[theorem]{Proposition}
\newtheorem{corollary}[theorem]{Corollary}

\theoremstyle{definition}
\newtheorem{definition}[theorem]{Definition}
\newtheorem{example}[theorem]{Example}

\theoremstyle{remark}
\newtheorem{remark}[theorem]{Remark}

%-------------------------
% Useful macros
%-------------------------
% \newcommand{\dd}{\mathrm{d}}                      % differential
\newcommand{\ii}{\mathrm{i}}                      % imaginary unit
\newcommand{\eexp}{\mathrm{e}}                    % exponential base
\DeclareMathOperator{\sgn}{sgn}

% Shortcuts for common physics notation (overrides safe if physics package present)
\renewcommand{\vec}[1]{\boldsymbol{#1}}

%-------------------------
% Bibliography (biber/biblatex) - optional
%-------------------------
% \usepackage[backend=biber,style=phys]{biblatex}
% \addbibresource{references.bib}

%-------------------------
% Document metadata
%-------------------------
\title{Execise 1}
\author{Zhuge X.K.}
\date{\today}

%-------------------------
% Document
%-------------------------
\begin{document}

\maketitle

Calculate the DOS for 2D \& 3D Dirac fermion using the Hamiltonians below:
\begin{equation*}
    \begin{aligned}
        H_{2D} &= \hbar v_F (k_x \sigma_x + k_y \sigma_y) \\
        H_{3D} &= \hbar v_F (k_x \sigma_x + k_y \sigma_y + k_z \sigma_z)
    \end{aligned}
\end{equation*}
\begin{proof}
    The DOS is defined as:
    \begin{equation*}
        D(E) = \int \frac{d^d k}{(2\pi)^d} \delta(E - E(k))
    \end{equation*}
    where $d$ is the dimension (2 or 3), and $E(k)$ is the eigenvalue of the Hamiltonian. For 2D case, the eigenvalue is:
    \begin{equation*}
        E = \pm \hbar v_F \sqrt{k_x^2 + k_y^2} = \pm \hbar v_F k
    \end{equation*}
    Thus,
    \begin{equation*}
        D(E) = \int \frac{d^2 k}{(2\pi)^2} \delta(E - \hbar v_F k) = \int_0^{2\pi} d\theta \int_0^\infty \frac{k dk}{(2\pi)^2} \delta(E - \hbar v_F k)
    \end{equation*}
    Using the property of delta function, we have:
    \begin{equation*}
        D(E) = \frac{1}{(2\pi)^2} 2\pi \frac{E}{\hbar^2 v_F^2} = \frac{E}{2\pi \hbar^2 v_F^2}
    \end{equation*}

    For 3D case, the eigenvalue is:
    \begin{equation*}
        E = \pm \hbar v_F \sqrt{k_x^2 + k_y^2 + k_z^2} = \pm \hbar v_F k
    \end{equation*}
    Thus,
    \begin{equation*}
        D(E) = \int \frac{d^3 k}{(2\pi)^3} \delta(E - \hbar v_F k) = \int_0^{\pi} d\theta \int_0^{2\pi} d\phi \int_0^\infty \frac{k^2 dk}{(2\pi)^3} \sin\theta \delta(E - \hbar v_F k)
    \end{equation*}
    Using the property of delta function, we have:
    \begin{equation*}
        D(E) = \frac{1}{(2\pi)^3} 4\pi \frac{E^2}{\hbar^3 v_F^3} = \frac{E^2}{2\pi^2 \hbar^3 v_F^3}
    \end{equation*}
\end{proof}

\end{document}